\section{Experimental Design: Pre-registered, Not Yet Run}
\label{sec:experiments}
We organize evaluation around \emph{falsification}, not around maximizing
a single benchmark score. No numerical outcomes for \spear{} are available
in this manuscript. The protocol and a fuller execution checklist are
provided in \code{EXPERIMENTS.md}.

\subsection{Research Questions and Testbeds}
\textbf{RQ1: Is external evidence causally useful for parameter updates?}
Can acquiring evidence improve the \emph{validated future-task utility}
of a parameter intervention over a trace-only update at matched cost?

\textbf{RQ2: Does parameter-aware source selection matter?}
When two failures have the same observable surface state, can a learned
policy select different evidence sources for different latent causes
and outperform fixed sources, confidence, and hand-written rules?

\textbf{RQ3: Is the improvement persistent and retention-aware?}
Does a committed update improve full trajectories on new task
instances with all auxiliary evidence disabled, without unacceptable
regression on previously learned skills?

\textbf{RQ4: Does the improvement process itself improve?}
With identical experience-log growth, does continuing to train
the query and intervention policies improve the cost to learn
\emph{later} skills compared with frozen meta-components?

\paragraph{Stage A: controlled diagnosis environment.}
Construct a resettable tool-use simulator with failures caused by
hidden schema changes, task-invariant violations, or environment-state
preconditions. Hold the visible error string and pre-query history
constant across selected pairs. Equip each source with distinct,
controlled access (manual snippet, advisory explanation, or executable
probe); test whether evidence changes the \emph{optimal update}
rather than merely identifying a label. Generate held-out
combinations of names, tool schemas, causal mechanisms, and
trajectory lengths. Randomize source availability/order to avoid
position shortcuts. A nontrivial subset must contain a
\emph{counterfactual reversal}: within the same visible failure
condition, the best validated update differs across hidden causes.

\paragraph{Stage B: interactive agent environments.}
Evaluate transfer to AppWorld \citep{trivedi2024appworld}, which
provides state-based unit tests for interactive tool use, and
AndroidWorld \citep{rawles2024androidworld}, which provides
parameterized mobile tasks and state-based checkers. Begin with
a compact instruction-tuned student (candidate:
Qwen3.5-4B, subject to reproducibility checks). Inject only
documented, reproducible capability gaps; distinguish benchmark
limitations (for example, failed resets or unavailable tool effects)
from agent failures. Use full-trajectory strict checker success,
not a language-model judgment, as the primary endpoint.
No outcomes from Stage A should be copied into a Stage B result.

\paragraph{Evaluation splits.}
Separate (i) intervention-proposal training, (ii) within-round
candidate-selection validation, (iii) meta-policy development,
and (iv) final untouched testing, at the \emph{mechanism/skill}
and task-instance levels. Keep old-skill retention suites fixed.
For sequential learning, reveal skill families in stages; later
test instances are unavailable during earlier adaptation.
Snapshot environments and version all evidence.

\subsection{Cost-Matched Comparators}
\label{sec:baselines}
\begin{table*}[t]
\centering
\small
\setlength{\tabcolsep}{5pt}
\begin{tabular}{p{3.3cm}p{4.2cm}p{4.0cm}p{4.2cm}}
\toprule
\textbf{Method} & \textbf{Evidence acquisition} &
\textbf{Parameter update} & \textbf{Isolated comparison}\\
\midrule
Unchanged student / memory-only &
None or experience retrieval &
None & Need for parameter change vs.\ external memory\\
Fixed Teacher + SFT &
Teacher on every eligible failure &
Gradient-based supervised update &
Standard verbal feedback and fine-tuning\\
Fixed Sources + Hypernetwork &
Fixed source order/bundle &
Evidence-conditioned adapter predictor &
Parameter generator without active evidence choice\\
Contextual BO / functional ES &
Budgeted trials selected from historical gain &
Low-dimensional update search &
Generic black-box optimization of updates\\
Heuristic / confidence routing &
Simple error rules or confidence thresholds &
Same update generator as full method &
Whether learned \pivoi{} beats trivial routing\\
Evidence Policy Only &
Learned source policy &
Frozen update generator &
Marginal benefit of acquisition learning\\
Intervention Model Only &
Fixed source policy &
Learned functional intervention generator &
Marginal benefit of intervention learning\\
Full \spear{} &
Learned parameter-aware source policy &
Learned functional generator; guarded updates &
Effect of joint decision coupling\\
Frozen-meta \spear{} &
Frozen policies; replay archive still grows &
Frozen generator; matched archive and budget &
Gain attributable to recursive learning\\
\bottomrule
\end{tabular}
\caption{Planned controls. All methods should receive a comparable
source-query, candidate-evaluation, and GPU-hour budget, and
the same versioned task stream. Results are intentionally absent.}
\label{tab:baselines}
\end{table*}

The most important comparisons are \emph{not} with the weakest
student: the main claim must survive (a) simple deterministic
routing, (b) fixed expert evidence, (c) a standard update mechanism
with an adaptive source selector, and (d) a strong learned update
mechanism with fixed sources. Where a named baseline has
unavailable code, describe and validate the faithful
reimplementation instead of claiming an exact reproduction.

\subsection{Metrics and Estimands}
\label{sec:metrics}
\begin{itemize}
\item \textbf{Primary endpoint:} strict end-to-end success on
      independent \emph{new} task instances after an accepted
      parameter update, at matched cumulative budget.
\item \textbf{Sample and cost efficiency:} verified skill gain per
      teacher call, retrieval call, simulator probe, GPU-hour,
      trained token, and total normalized expense.
\item \textbf{Retention:} average drop and worst-skill drop on the
      pre-registered old-task suite, with side-effect violations.
\item \textbf{Value prediction:} signed prediction error, ranking
      regret, calibration, and false-acceptance rate of
      $Q_\omega$ and estimated \pivoi{}.
\item \textbf{Source-policy increment:} paired utility relative to
      fixed sources, confidence, and best hand-written rule
      under matched queries and candidate counts.
\item \textbf{Recursive gain:} area under the skill-acquisition
      curve, later-skill improvement per expense, and adaptive
      minus frozen-meta learning efficiency at identical replay size.
\end{itemize}

The primary causal estimand is the \emph{difference in post-update
utility under a randomized evidence-source policy} at a fixed
checkpoint and intervention-validation budget. Report results by
latent cause and by visible-failure stratum, not just in aggregate.
A source-selection effect should include a check that update choices
actually differ after evidence; success alone may be due to more
computation.

\subsection{Paired Evaluation and Statistical Reporting}
At each eligible decision, branch from the same checkpoint and
resettable environment. Use paired seeds for untreated and candidate
updates, randomized source assignments where valid, and
independent final evaluation seeds. Publish the full denominator,
technical-error exclusions, source propensities, cost ledger,
and confidence intervals. A priori sample size should be chosen
through pilot-based power analysis; the final test set and
analysis rules must be frozen before model selection.
Use a hierarchical or paired bootstrap across task mechanisms
and random seeds, with a pre-specified correction for multiple
primary comparisons. Report inconclusive results as inconclusive.

\subsection{Falsification Gates and Critical Ablations}
The proposed sequence includes four \emph{go/no-go gates}:
\begin{enumerate}
\item \textbf{Evidence gate.} At fixed update budget, new external
      evidence improves validated adaptation over trace-only
      updates; distinguish this from paraphrasing the trace.
\item \textbf{Conditional reversal gate.} Within matched observable
      conditions, the best intervention (including no update)
      genuinely changes with additional evidence. If a
      one-condition rule is sufficient, do not train a large router.
\item \textbf{Persistence gate.} A committed update retains gains on
      new full trajectories with consultant text and retrieval off,
      and does not breach the old-skill regression tolerance.
\item \textbf{Recursion gate.} Learning $\eta$ and $\phi$ improves
      later update quality beyond the same replay archive with
      both components frozen.
\end{enumerate}

Run text-shuffled, source-masked, trace-only paraphrase,
oracle-evidence, basis-only, residual-only, no-update, and
no-retention-gate ablations. Randomize nominal tool names and
document layouts; hold out mechanism combinations so gains
cannot be attributed to a memorized lookup table. If the gates
fail, the corresponding novelty or RSI claim must be narrowed.

\paragraph{Result-reporting policy.}
The evaluation tables above specify \emph{methods and outcomes to
measure}; they contain no experimental values. A future Results
section should be added only with artifact-linked raw runs,
strict-checker audit trails, and the declared statistical analysis.

\paragraph{Executable mechanism check.}
A separately marked, CPU-only prototype of the controlled
counterfactual protocol, including real-run audit schemas, is available
under \code{experiments/parameter\_voi/}. This prototype uses a
synthetic two-weight student and simulated evidence sources: it does
\emph{not} constitute an experiment on LLM agents or evidence that
any of the paper-grade gates above have been met.
